%
% OBJECTS
%
\chapter[Espace, temps, lois]{Vers des objets cellulaires:~espace, temps, lois}
\mysetref{space-time-law}

Nous soutenons qu'un environnement de développement d'AC doit reposer
sur un ensemble d'objets permettant la construction d'expériences. %
Ces objets doivent correspondre aux ressources utilisables par une
simulation et fournir les opérations et les fonctions de réduction
d'information utilisées lors de la description de l'expérience
projetée. %

Les objets de base d'une simulation sont l'\aconcept{espace}, %
le \aconcept{temps} et la \aconcept{fonc\-tion de transition}. %
La fonction de transition se compose d'une structure d'interconnexion
cellulaire associable à un espace, d'une fonction de transition
locale\footnote{%
  L'expression \aquote{fonction de transition} désigne, suivant le
  contexte, soit l'objet regroupant le réseau d'interconnexion et la
  fonction de transition locale, soit la fonction de transition locale
  elle-même. %
  La fonction de transition locale sera par ailleurs également nommée
  \aquote{règle} ou \aquote{table de transition}, et l'objet fonction
  de transition sera à l'occasion nommé loi, ou AC.}, %
applicable de manière synchrone sur l'ensemble des cellules d'un
espace doté d'un réseau d'interconnexion, et d'une fonction
d'application utilisant les composants de la
fonction de transition pour générer le nouvel état d'un l'espace. %
La fonction de transition réalise les lois de la dynamique de
l'univers cellulaire simulé. %
Elle est génératrice du temps par application sur l'espace. %

Des instances d'objets de base sont regroupées pour la réalisation
d'une simulation au sein d'un objet de type séquenceur. %
Des manipulateurs et des observateurs peuvent être associés aux objets
de la simulation. %

%%
%% OBJECTS/SPACE
%%
\section{L'espace:~la mémoire}

L'espace est un ensemble de cellules ordonnées conformément à la
structure d'interconnexion de la fonction. %
La structure de l'espace peut théoriquement être quelconque pourvu
qu'elle soit homogène. %
En pratique nous limitons initialement les structures disponibles aux
espaces de dimension~$d$~1, 2~ou~3. %
\mysetref{nombre-de-voisin}%
Le nombre~$r$ de cellules distantes de $n$~unités au plus d'une
cellule quelconque de l'espace est $(2n + 1)^{d}$. %
La structure d'interconnexion sera un sous-ensemble de l'ensemble des
cellules distantes de $n$~unités au plus, avec~$n$ typiquement limité
à 1~ou~2 unités. %
Nous appellerons \anewterm{cube d'interconnexion} de dimension~$d$ une
structure d'interconnexion complète pour la dimension~$d$ et
\anewterm{croix d'interconnexion} de dimension~$d$ l'ensemble des
cellules distantes de $n$~unités au plus sur une dimension de l'espace
au plus. %
Le cube d'interconnexion de dimension~2 et distance~1 correspond au
voisinage de Moore, et la croix d'interconnexion de dimension~2 et
distance~1 au voisinage de von Neumann. %
Le nombre d'éléments d'une croix d'interconnexion est $2 n d + 1$.
%%
\mytable{neighburhood-size-cube}%
{Nombres d'états pour quelques cubes d'interconnexions,}%
{Tailles induites par des cubes d'interconnexions}%
{Taille en octets des tables de transitions de quelques cubes
  d'interconnexions de cellules binaires.}

Le tableau~\mygettbl{neighburhood-size-cube} donne la taille en bytes
des tables de transitions des cubes d'interconnexions de cellules
binaires pour les 4~premières dimensions d'espace et les 4~premières
distances. %
Les entrées en caractères gras sont réalisables par lookup-table.

Le tableau \mygettbl{neighburhood-size-croix} suggère l'intérêt de
voisinages mixtes
\[\text{croix}(d, n) \cup \text{cube}(d, n - 1)\]

Mais la structure d'interconnexion régulière offrant le meilleur
compromis \trialt{distance}{volume}{isotropie} est une \anewterm{sphère de
  Manhattan}, une hyper-sphère de dimension~$d$ dont le rayon est une
distance de Manhattan\footnote{%
  Le volume d'une sphère de Manhattan est %
  $S(d, r) = S(d - 1, r) + 2 \sum_{i=0}^{i=R-1} S(d - 1, i)$ %
  avec $S(1, r) = 2r + 1$. %
  Les définitions que nous proposons pour les cubes, croix et sphère de
  connexions considèrent le centre comme exclu du rayon. %
  L'inclusion du centre dans le rayon permet la définition de cubes et
  de sphères privés d'un centre ponctuel. %
  Le voisinage de Margolus~\cite[]{Toffoli87} par exemple utilise un
  carré de 2~cellules de côté.}. %
Une sphère de Manhattan de rayon~1 définit la même structure
d'interconnexion qu'une croix de rayon~1. %
Le voisinage de von Neumann est donc une sphère de Manhattan, mais pas
un voisinage de Moore. %
Avec des cellules binaires, en dimension~2 les sphères de rayon
2~et~3 sont accessibles\footnote{1K et 4M}, en dimension~3 la sphère
de rayon~2 est accessible\footnote{4M}.

\noindent L'espace supporte des fonctions d'\bialt{entrée}{sortie},
d'édition, de réduction et de transformation.
%%
\mytable{neighburhood-size-croix}%
{quelques croix d'interconnexions,}%
{Tailles induites par des croix d'interconnexions}%
{La taille en bytes des tables de transitions induite par des croix
  d'interconnexions de cellules binaires permet l'exploration de
  voisinages utilisant de plus grandes valeurs pour les paramètres
  dimension et rayon, au détriment du rapport
  $\frac{\text{volume}}{\text{rayon}}$, donc de l'isotropie.}
%%
\mytable{neighburhood-size-sphere}%
{et quelques sphères d'interconnexions.}%
{Tailles induites par des sphères d'interconnexions}%
{La taille en bytes des tables de transitions induite par des sphères
  d'interconnexions de cellules binaires, est particulièrement
  intéressante pour les valeurs $d=2, n=3$ et $d=3, n=2$ qui sont
  inaccessibles au voisinage cube \mygettbl[paren]{neighburhood-size-cube}.}
%%

%%
%% OBJECTS/SPACE/STRUCT
%%
\subsection{Structure de l'espace:~dimension et voisinage}

L'espace apparaît comme une ligne, un plan, ou un volume de cellules
adres\-sables globalement par leurs coordonnées euclidiennes. %
Nous utiliserons un espace de dimension~2 par défaut pour exposer les
opérations applicables sur les espaces. %
Cet espace peut être vu et manipulé soit comme un plan de cellules,
soit comme un vecteur de plans de bits correspondant à une tranche
cellulaire. %
Les opérations sur un espace sont toutes les opérations possibles sur
une mémoire organisée selon une dualité vecteur de plans de bits, plan
de vecteurs de bits.

Idéalement, en association avec une structure d'interconnexion, le
vecteur de bits de la mémoire cellulaire est associé à une structure
de nommage. %
Les noms des zones de mémoire sont utilisables en paramètres des
opérations portant sur les tranches d'espaces cellulaires~(édition,
visualisation, réduction) et sur les variables cellulaires~(définition
de la fonction de transition).

%%
%% OBJECTS/SPACE/EDIT
%%
\subsection{Édition d'espace:~les configurations cellulaires}

Une première famille d'opérations est fournie par un éditeur de
\astrangeterm{pixmaps}{%
  Une mémoire de pixels.} %
permettant la construction des configurations d'espaces
au temps zéro d'une simulation. %
L'éditeur permet également la modification du point d'espace courant
dans le temps entre chaque application de la fonction de transition.
%%
%% neighburhood-size-tbl-16xor9
%%
\mydumpfullR[htb]{.75}{neighburhood-size-tbl-16xor9}%
{Variété des configurations initiales}%
{Édition d'espace;~génération~16 de la règle \arule{Xor9} sur le
  tableau~\mygettbl{neighburhood-size-cube}}%
{Nous avons choisi \mysoftcite[PBM]{pbm} comme format\footnote{%
    Les numéros en indice supérieur entre crochets renvoient au
    répertoire des outils~\mygetref[paren]{tools}.} %
  pivot de \bialt{bitmap}{pixmap}. %
  Cette figure montre l'application de la règle~\myrulecite{xor9} sur
  un espace obtenu à partir du tableau~\mygettbl{neighburhood-size-cube}. %
  Pour ce faire \mysoftcite{groff} est utilisé pour produire une
  représentation PostScript de la page, puis \mysoftcite{ghostscript}
  est utilisé pour produire une représentation \asoft{PBM} de la page,
  divers outils comprenant le format \asoft{PBM} tel \mysoftcite{xv}
  peuvent être utilisés à ce stade pour retravailler le bitmap, enfin
  le bitmap est utilisé par le workbench qui fournit le bitmap
  transformé par l'application itérative d'une règle, lequel est
  retransformé en PostScript pour insertion dans le document final.}
%%
Il est souhaitable de pouvoir utiliser aussi d'autres outils que les
nôtres~(que nous nommons outils externes) pour manipuler les espaces. %
Le format de stockage externe doit permettre l'application
\astrangeterm{off-line}{Le terme off-line désigne le passage des
  objets de la simulation à d'autres processus.} de nombreux outils
sur les espaces. %
Ceci permet une intégration étroite de ces outils à l'environnement de
développement \bialt{et}{ou} d'interaction, mais implique l'existence de
plusieurs classes de débit pour les opérations applicables sur les
objets. %
Cette contrainte est quasiment sans effet pour l'utilisation
d'éditeurs interactifs de type \astrangeterm{painter}{%
  Un outil métaphore de la palette et du pinceau, lorsque ceux-ci sont utilisés
  pour peindre une surface} %
ou \astrangeterm{drawer}{%
  Un outil métaphore du crayon
  lorsqu'il est utilisé pour dessiner.} %
sur un espace.
%%
L'éditeur devrait également fournir une famille de fonctions de
remplissage aléatoire de l'espace aussi riche que possible.
%%

%%
%% OBJECTS/SPACE/REDUCT
%%
\subsection{Réduction d'espace:~taille et contenu}

La deuxième famille d'opérations fournit des fonctions de réduction de
l'information contenu dans un espace. %
Il s'agit de réductions arithmétiques élémentaires~(somme,
histogramme, parité) de réductions statistiques standard~(minimum,
maximum, moyenne, entropie, dimension fractale) ou spécialisées comme
le comptage de patterns propres à une simulation. %
On peut ajouter les outils standards de compressions de données qui,
outre l'archivage des espaces, fournissent des quantités
informationnelles.

\noindent Ici encore l'utilisation d'un format pivot standard et doté
de nombreux outils de conversion de format permet l'utilisation
d'outils statistiques off-line.\footnote{%
  \mygettbl[phrase]{lifex-sum-1}} %
De plus les outils externes se présentant sous la forme de librairie
sont également intégrables à l'environnement de développement et
permettent une utilisation directe qui ne nécessite pas l'emploi de
fichiers temporaires.

%%
%% OBJECTS/SPACE/TRANSFORM
%%
\subsection{Transformation d'espace:~d'une forme à l'autre}

%%
%% tsum-128-edge
%%
\mydumpfull{tsum-128-edge}%
{Transformation d'espace}%
{Transformation d'espace;~un dithering et deux détections de contours}%
{L'espace original est présenté par la figure \mygetdump{tsum-grey}. %
  Ces transformations ont été réalisées avec les outils \asoft{PBM}.
  Le \aquote{dithering} permet la réduction d'un greymap à un bitmap.}
%%
Enfin, des transformations globales arbitraires sont applicables sur
les espaces. %
Il s'agit par exemple des fonctions de traitement du signal en
général, et de traitement d'images tel que la reconnaissance des
formes et l'extraction de contour. %
La figure \mygetdump{tsum-128-edge} montre une détection de contours
utilisant des outils fournis par \asoft{PBM}. %
\mygetdump[seealso]{tsum-grey}

Certaines de ces fonctions sont elles mêmes basées sur des algorithmes
cellulaires, elle doivent alors pouvoir être réalisées à l'aide de
l'environnement de développement, les fonctions de séquencement
permettant leur application non destructive à la volée sur le point
d'espace courant d'une séquence.

%%
%% OBJECTS/SPACE/VISUAL
%%
\subsection{Visualisation d'espace:~comprendre les lois}
\mysetref{obj-visualisation}

L'espace constitue un objet de visualisation élémentaire a priori. %
Il peut être vu par l'intermédiaire de fenêtres X11 associées à un
bitmap, ou à une séquence de plans de bits~(pixmap) associée à une
\aquote{colormap}.
%%
\mydumpfullR[ht]{.75}{anneal-g128-s11}%
{Visualisations par bitmap:~d'un espace,}%
{Un bitmap point-bit~\asquare{2048}}%
{Génération~128 de la règle~\arule{Anneal}, sur une espace initial de
  densité\,\Sundemi, et de taille~\asquare{2048}.}
%%

\noindent%
La figure~\mygetdump{anneal-g128-s11} montre la visualisation par
bitmap de l'espace de la géné\-ration~128 de la
règle~\arule{Anneal}\footnote{%
  \mygetrule[phrase]{anneal}}. %
L'espace est de taille~\asquare{2048}. %
L'espace initial a une densité aléatoire\,\undemi. %
La taille de la cellule étant de 1~bit, le bitmap permet de visualiser
l'intégralité des états cellulaires. %
Chaque point de l'image correspond à un bit.

\noindent%
La figure~\mygetdump{tube-worm-mix} montre 8~séquences de
8~espaces chacune. %
Il s'agit de séquences de la règle~\arule{Tube-Worm}\footnote{%
  \mygetrule[phrase]{tube-worm}} %
appliquée sur le même espace initial, les variations portant sur les
paramètres statiques\footnote{%
  \mygetref[phrase]{static-param}} %
de la règle. %
Le nombre de générations entre deux espaces dépend d'un des paramètres
statiques de la règle. %
Il est de 32~intervalles du temps propre à cette règle.

\noindent%
La figure~\mygetdump{tsum-grey} montre la visualisation par pixmap de
2~générations de la règle \arule{Qmean} \mygetrule[paren]{qmean}. %
La taille de la cellule étant de 8~bits, 256~couleurs sont nécessaires
pour visualiser tous les états possibles\footnote{%
  Bien entendu, en termes purement informationnel l'espace est comme
  pour la  figure~\mygetdump{anneal-g128-s11}~\asquare{2048}.}.

\noindent%
La figure \mygetdump{tsum} montre la visualisation des
générations~$2^{3}$ à~$2^{10}$~(de haut en bas) de la
règle~\arule{Qmean}. %
L'espace est de taille \asquare{128}. %
L'espace initial a une densité aléatoire\,\undemi. %
Les 256~états possibles de chaque cellule sont visualisés par
8~bitmaps au lieu d'un pixmap. %
Les bitmaps sont présentés de gauche à droite pour les bits~0 à~7.

Les variables de réductions d'espace sont présentées sous forme ASCII,
associées aux fenêtres pixmap, ou éventuellement par des
\astrangeterm{viewers}{%
  visualiseurs.} %
de groupe de variables scalaires~(bargraph par exemple). %
Ce type de visualisation~(bitmaps + scalaires) est assez rapide pour
permettre la réalisation de films de séquences d'espaces à une
fréquence supérieure à 1~Hz.

D'autres types de visualisations dédiées à une simulation particulière
ou plus consommatrice en ressources sont possibles:
\begin{itemize}
\item Une version iconique de l'état d'une cellule.
\item Une vision tridimensionnelle, réalisée par \astrangeterm{plotter}%
  {Traceur. Ce type de visualisation repose évidemment sur le
    transfert d'un espace à un outil externe spécialisé. %
    Nous avons utilisé \mysoftcite{gnuplot} et \mysoftcite{Khoros}.} %
  \mygetplot[paren]{anneal-sum} ou \astrangeterm{renderer}%
  {Générateur de rendu. Nous avons utilisé \mysoftcite{qrt} et la toolbox
    \mysoftcite{vart} sous \asoft{Khoros}.} %
  \mygetdump[paren]{tsum-ray-one} d'un espace considéré comme surface
  d'élévation, ou dont une des variables d'état de la cellule est
  considérée comme donnée d'élévation, et une autre comme donnée de
  projection.
\item Pour les espaces de dimension~3, une vision tridimensionnelle
  par seuillage, épluchage ou transparence.
\end{itemize}
Les visualisations coûteuses en ressources, ne permettant pas la
réalisation de films temps réel, servent à la réalisation
d'instantanés, la réalisation de films de séquences d'instantanés
restant possibles en temps différé.
%%
\mydumpfull{tube-worm-mix}%
{et de plusieurs séquences d'espaces.}%
{Un bitmap constitué de $64$~bitmaps~\asquare{256}}%
{Une composition construite avec 8~séquences~(de haut en bas), de
  8~espaces chacune~(de gauche à droite) de la règle
\arule{Tube-Worm} appliquée sur le même espace initial}
%%
\mydumpfull{tsum-grey}%
{Visualisations d'espaces par un pixmap}
{Visualisation par pixmap de 2~générations de la règle \arule{Qmean}}%
{Les génération~128 et~256~(de gauche à droite) de la règle
  \arule{Qmean} utilisant deux types de colormaps~(de haut en bas). %
  L'espace est de taille~\asquare{256}. %
  L'espace initial a une densité aléatoire\,\Sundemi. %
  Les colormaps sont ici des greymaps. %
  Les deux pixmaps supérieurs utilisent une greymap dont l'échelle des
  valeurs est progressive. %
  Les deux pixmaps inférieurs utilisent une greymap dont l'échelle des
  valeurs est aléatoire, masquant la répartition générale des
  intensités, mais révélant les détails de l'alternance
  de valeurs.}
%%
\mydumpfull{tsum}%
{ou par plusieurs bitmaps}%
{Visualisation par bitmap de 8~plans de 8~générations de la
  règle~\arule{Qmean}}%
{Les générations~$2^{3}$ à~$2^{10}$~(de haut en bas) de la règle
  \arule{Qmean}. %
  L'espace est de taille~\asquare{128}. %
  L'espace initial a une densité aléatoire\,\Sundemi. %
  Les 256~états possibles de chaque cellule sont visualisés par
  8~bitmaps au lieu d'un pixmap. %
  Les bitmaps sont présentés de gauche à droite pour les bits~0 à~7.}
%%
\mydumpfull{tsum-ray-one}%
{Visualisation par lancer de rayons~(ray tracing)}
{Visualisation par lancé de rayons~(ray tracing) d'espaces de la
  règle \arule{Qmean}}%
{La fonction d'extraction d'isosurface fournie par \mysoftcite{vart} a
  été utilisée sur le volume constitué par l'empilement de
  128~espaces~(128~générations) de~\asquare{128} cellules. %
  L'espace initial est issu de la génération~8192 de la
  règle~\arule{Anneal} elle-même appliquée sur un espace de
  taille~\asquare{1024} et de densité aléatoire\,\Sundemi, réduit
  à~\asquare{128}.}
%% 
\mydumpfullR{.8}{tsum-ray-all}%
{Visualisation par films de vues ray-tracées}%
{Visualisation par films de vues ray-tracées d'espaces de la
  règle \arule{Qmean}}%
{Par rapport à la figure \mygetdump{tsum-ray-one}, les variations de
  paramètres sur l'axe horizontal porte sur la prise de vues; %
  les variations de paramètres sur l'axe vertical porte sur la valeur
  de seuillage utilisée. %
  La valeur de seuil varie de 0.3~à~0.8.}
%%
\clearpage

%%
%% OBJECTS/SPACE/RESSOURCE
%%
\subsection{Ressources associées à l'espace:~des limites à choisir}

La consommation mémoire d'un espace d'arête~$a$ et de dimension~$d$
est~$a^{d}$. %
Pour des raisons détaillées lors de la présentation de
la fonction de tran\-si\-tion\footnote{%
 \mygetref[phrase]{mem-struct}}, %
l'organisation mémoire de l'espace privilégie le plan de vecteurs de
bits, d'une taille multiple de l'unité adressable, c'est à dire que
les cellules seront constituées au minimum d'un byte.
%%
\mytable{space-size}
{Taille d'espaces fonction de la dimension et
  de l'arête}%
{Tailles d'espaces pour quelques tailles d'arête des premières dimensions}%
{Ce tableau donne les tailles en byte, fonction des paramètres dimension
  et arête, d'espaces accessibles à l'expérimentation.}

Le tableau~\mygettbl{space-size} donne un aperçu des tailles d'espaces
maximum accessibles pour les 4~premières dimensions d'espace. %
On peut voir que si les espaces de dimension~2 ne posent pas de
problèmes, les espaces de dimension~3 seront avantageusement ramenés à
des parties de cube. %
Par exemple, les~256K d'un espace~$128 \times  64 \times  32$ sont probablement
plus avantageux que les~256K d'un espace~$64 \times  64 \times  64$. %
La consommation CPU de la fonction de transition, croissant
linéairement en fonction de la taille de l'espace, on peut également
avoir une idée des temps de calculs associés à la taille des espaces. %
Celui-ci est $\frac{\text{G} a^{d}}{\text{D}}$, avec $\text{G}$ pour
le nombre de générations, $d$ pour la dimension de l'espace et
$\text{D}$ pour le débit de la fonction d'application.
%%
\mytable{compute-time-100kcps}%
{Temps de calcul sur ces espaces de 1000~générations à 100~KCPS,}%
{Temps de calcul sur les espaces du tableau~\mygettbl{space-size}
  fonction du nombre de génération et de la puissance de calcul}%
{Ce tableau donne un aperçu des expériences réalisables fonction des
  tailles d'espace et de la puissance de calcul du moteur cellulaire. %
  Les paramètres nombre de générations et puissance de calcul sont
  fixés respectivement à 1000~générations et 100~KCPS\@. %
  Un temps de calcul d'une journée pour 1000~générations est considéré
  comme la limite praticable. %
  Un débit à 100~KCPS constitue une hypothèse basse. %
  Notre environnement fournit actuellement un débit supérieur d'un
  ordre de grandeur.}
%%
\mytable{compute-time-10mcps}%
{et de 1000~générations à 10~MCPS.}%
{Variations des temps de calculs du
  tableau~\mygettbl{compute-time-100kcps} résultant d'une augmentation
  de deux ordre de grandeur de la puissance de calcul}%
{Le gain d'un ordre de grandeur, par rapport au performance de notre
  implémentation, ouvrirait l'exploration des espaces de
  dimension~3~(nous considérons empiriquement qu'une arête de
  taille~256 constitue une base suffisante à l'exploration de ces
  espaces). %
  Ce gain pourrait être obtenu en partie par une
  implémentation répartie du moteur cellulaire.}

Le tableau~\mygettbl{compute-time-100kcps} montre les temps associés
au calcul de 1K\footnote{%
  Les débits sont en puissance
  de~10:\hfil\break 1~kilocell = 1024~cell, %
  1~kilo~génération = 1000~générations.}%
~générations d'espaces du tableau \mygettbl{space-size} avec une
fonction d'application à 100~KCPS\footnote{%
  Kilo Cell update Per
  Second.}, (consommant 100~instructions par cellule sur une machine
10~MIPS, par exemple). %
Seule l'unité de temps la plus significative est indiquée. %
On pourrait obtenir un débit de 10~MCPS avec une fonction
d'application consommant 20~instructions par cellule, et quatre
processeurs à 50~MIPS, lesquels ne partagent qu'une faible partie de
la mémoire totale\footnote{%
  \mygetref[phrase]{repartition-du-calcul}}. %
Le tableau \mygettbl{compute-time-10mcps} reprend les calculs du
tableau \mygettbl{compute-time-10mcps} pour %
$\text{G} = 10 \text{MCPS}$.
%%
\clearpage

%%
%% Local Variables:
%% mode: latex
%% mode: auto-fill
%% iso-accents-minor-mode: nil
%% TeX-master: "top"
%% TeX-command-default: "mytex"
%% Local IspellParsing: "latex-mode ~latin1"
%% Local IspellPersDict: "~/.ispell_lisp"
%% LocalWords: "Moore von Neumann neighburhood-size bytes lookup isotropie xor"
%% LocalWords: "Manhattan hyper-sph\`ere Margolus adressables euclidiennes nommage"
%% LocalWords: "pixmaps off-line off-line painter drawer neighburhood-size-tbl"
%% LocalWords: "PBM pbm bitmap pixmap rule groff groff ghostscript ghostscript"
%% LocalWords: "xv xv CAW fractale patterns informationnelles tbl lifex-sum tsum"
%% LocalWords: "int\'egrables edge dithering dump tsum-grey s\'equencement colormap"
%% LocalWords: "destructive anneal-g point-bit anneal anneal tube-worm-mix qmean"
%% LocalWords: "bitmaps textbf tube-worm static-param greymaps greymap viewers"
%% LocalWords: "informationnel bargraph iconique plotter gnuplot gnuplot Khoros"
%% LocalWords: "khoros anneal-sum renderer qrt qrt toolbox Vart vart seuillage"
%% LocalWords: "tsum-ray-one ray-tracer tsum-ray-all ray-trac\'ees d'isosurface cu"
%% LocalWords: "mem-struct adressable byte space-size compute-time kcps per ht"
%% LocalWords: "KCPS cell cell update MIPS MCPS mcps nombre-de-voisin paren"
%% LocalWords: "neighburhood-size-cube lookup-table neighburhood-size-croix ray"
%% LocalWords: "neighburhood-size-sphere tools workbench seealso visualiseurs"
%% LocalWords: "Traceur tracing impl\'ementation kilocell repartition-du-calcul"
%% End:
